\begin{definition}[Inductive Subset of a Peano System]
\label{def:inductive-subset-of-peano-system}
Let $(P,S,1)$ be a Peano system, and let $A \subseteq P$. The subset $A$ is
called \emph{inductive} exactly when $1 \in A$ and $A$ is successor-closed.
\end{definition}
\end{definitionbox}
\LeanFormalizes{def:inductive-subset-of-peano-system}{lra-lean}{LRA.VolumeII.PeanoSystems.TexSkeleton}{InductiveSubsetOfPeanoSystem}{pending}

\begin{remark*}[Standard quantified statement]
For a Peano system $(P,S,1)$ and a subset $A\subseteq P$,
\[
\bigl(1 \in A \land
\forall n \in A,\ S(n) \in A\bigr).
\]
\end{remark*}

\begin{remark*}[Negated quantified statement]
\[
\begin{aligned}
&\text{Let } (P,S,1) \text{ be a Peano system and let } A \subseteq P.\\
&A \text{ is not inductive}
\Longleftrightarrow
\bigl(1 \notin A \lor \exists n \in A,\ S(n) \notin A\bigr).
\end{aligned}
\]
\end{remark*}

\begin{remark*}[Failure modes]
\begin{description}
  \item[Exposition.]
  A subset fails to be inductive in exactly two ways: it may fail to contain
  the distinguished first element, or it may contain some element while
  failing to contain that element's successor.
  \item[Quantified failure.]
  \[
  \begin{aligned}
  A \text{ is not inductive}
  \Longleftrightarrow\;&
  \underbrace{1 \notin A}_{\text{base point missing}}\\
  &\lor
  \underbrace{\exists n \in A,\ S(n) \notin A}_{\text{successor closure fails}}.
  \end{aligned}
  \]
\end{description}
\end{remark*}

\begin{remark*}[Interpretation]
An inductive subset contains the distinguished first element and is closed
under successor. Such a subset contains every stage obtained by starting at
$1$ and applying successor repeatedly.
\end{remark*}

\begin{remark*}[Source comparison]
The notion of an inductive subset packages the two hypotheses appearing in the
induction clause of Feferman's Definition 3.1. The present notes name this
package so that the internal structure of induction proofs remains explicit.
\citep{FefermanNumberSystems1964}
\end{remark*}

\begin{dependencies}
\begin{itemize}
  \item \hyperref[def:peano-system]{Peano System}
  \item \hyperref[def:successor-closed-subset]{Successor-Closed Subset}
\end{itemize}
\end{dependencies}

